% GENERATED FILE. Edit canonical lessons, then run npm run generate:books.
% canonical-source-hash: fnv1a64:844813320baaac64
% canonical-lessons: MW-C207-bagicho, MW-C207-kund, MW-C207-nako, MW-C207-phatak, MW-C207-buraj

\chapter{A Garden and a Tank}
\label{ch:mw-a-garden-and-a-tank}
\hlchaptermodality{207}

\begin{chapteropening}
\textbf{By the end of this chapter:} \emph{I can name the places around the old town.}

Complete the last lesson of chapter 207: I can name the places around the old town.
\end{chapteropening}

\section[\texorpdfstring{\mw{बगीचो} (bagīcho)}{बगीचो (bagīcho)}]{\mw{बगीचो} (bagīcho) — a garden}

\label{lesson:MW-C207-bagicho}



\begin{quote}
Before the new one: say the Marwadi for on the other side, then the Marwadi for separate.
\end{quote}

\subsection*{You'll want to know: \mw{बगीचो}}
\textbf{\mw{बगीचो}} — \emph{bagīcho} — \textquotedblleft{}a garden\textquotedblright{}.

\begin{quote}
\mw{बगीचो} \mw{कोट} \mw{कनै} \mw{है।}
\end{quote}

\begin{quote}
\emph{bagīcho koṭ kanai hai.}
\end{quote}

\begin{quote}
The garden is near the fort.
\end{quote}

A masculine noun.

\begin{cousinweb}
\textbf{\mw{बगीचो}} (\emph{bagīcho}) is usually traced to Persian \emph{bāghcha}, \textquotedblleft{}a small garden\textquotedblright{}.
\end{cousinweb}

\begin{grammarlens}[title={What you've built}]
One more word for this chapter.
\end{grammarlens}

\subsection*{Your turn}
\begin{itemize}
\raggedright
  \item \emph{Say it:} \emph{bagīcho}
  \item \emph{Say it:} \emph{bagīcho}, once more
  \item \emph{Say it:} \emph{bagīcho koṭ kanai hai.} — The garden is near the fort.
  \item \emph{Recall:} say the Marwadi for on the other side, then the Marwadi for separate, then say \emph{bagīcho} again
\end{itemize}

\begin{tcolorbox}[breakable,colback=teal!4,colframe=teal!35!black,title={Before you move on}]
What does \mw{बगीचो} mean? (\textquotedblleft{}A garden\textquotedblright{}.) Say it once more.
\end{tcolorbox}
\section[\texorpdfstring{\mw{कुंड} (kuṇḍ)}{कुंड (kuṇḍ)}]{\mw{कुंड} (kuṇḍ) — a water tank}

\label{lesson:MW-C207-kund}



\begin{quote}
Before the new one: say the Marwadi for separate, then the Marwadi for a garden.
\end{quote}

\subsection*{You'll want to know: \mw{कुंड}}
\textbf{\mw{कुंड}} — \emph{kuṇḍ} — \textquotedblleft{}a water tank\textquotedblright{}.

\begin{quote}
\mw{कुंड} \mw{मांय} \mw{पाणी} \mw{कोनी।}
\end{quote}

\begin{quote}
\emph{kuṇḍ mā̃y pāṇī konī.}
\end{quote}

\begin{quote}
There's no water in the tank.
\end{quote}

A masculine noun.

\begin{cousinweb}
\textbf{\mw{कुंड}} (\emph{kuṇḍ}) is inherited from Sanskrit \emph{kuṇḍa}, \textquotedblleft{}bowl, basin of water\textquotedblright{}. CDIAL gives no Marwari form in this entry, but lists Hindi \emph{kū̃ḍ}, \textquotedblleft{}tub\textquotedblright{}, and Gujarati \emph{kũḍ}, \textquotedblleft{}basin\textquotedblright{} (CDIAL 3264).
\end{cousinweb}

\begin{grammarlens}[title={What you've built}]
One more word for this chapter.
\end{grammarlens}

\subsection*{Your turn}
\begin{itemize}
\raggedright
  \item \emph{Say it:} \emph{kuṇḍ}
  \item \emph{Say it:} \emph{kuṇḍ}, once more
  \item \emph{Say it:} \emph{kuṇḍ mā̃y pāṇī konī.} — There's no water in the tank.
  \item \emph{Recall:} say the Marwadi for separate, then the Marwadi for a garden, then say \emph{kuṇḍ} again
\end{itemize}

\begin{tcolorbox}[breakable,colback=teal!4,colframe=teal!35!black,title={Before you move on}]
What does \mw{कुंड} mean? (\textquotedblleft{}A water tank\textquotedblright{}.) Say it once more.
\end{tcolorbox}
\section[\texorpdfstring{\mw{नाको} (nāko)}{नाको (nāko)}]{\mw{नाको} (nāko) — a checkpoint; the mouth of a road}

\label{lesson:MW-C207-nako}



\begin{quote}
Before the new one: say the Marwadi for a garden, then the Marwadi for a water tank.
\end{quote}

\subsection*{You'll want to know: \mw{नाको}}
\textbf{\mw{नाको}} — \emph{nāko} — \textquotedblleft{}a checkpoint; the mouth of a road\textquotedblright{}.

\begin{quote}
\mw{नाको} \mw{अठै} \mw{सूं} \mw{नजीक} \mw{है।}
\end{quote}

\begin{quote}
\emph{nāko aṭhai sū̃ najīk hai.}
\end{quote}

\begin{quote}
The checkpoint is close to here.
\end{quote}

A masculine noun.

\begin{cousinweb}
\textbf{\mw{नाको}} (\emph{nāko}) is inherited from a reconstructed \emph{*nākka}, \textquotedblleft{}like a nose, projecting\textquotedblright{}. CDIAL gives no Marwari form in this entry, but lists Hindi \emph{nākā}, \textquotedblleft{}end of road, passage\textquotedblright{}, and Sindhi \emph{nāko}, \textquotedblleft{}entrance, pass\textquotedblright{} (CDIAL 7037).
\end{cousinweb}

\begin{grammarlens}[title={What you've built}]
One more word for this chapter.
\end{grammarlens}

\subsection*{Your turn}
\begin{itemize}
\raggedright
  \item \emph{Say it:} \emph{nāko}
  \item \emph{Say it:} \emph{nāko}, once more
  \item \emph{Say it:} \emph{nāko aṭhai sū̃ najīk hai.} — The checkpoint is close to here.
  \item \emph{Recall:} say the Marwadi for a garden, then the Marwadi for a water tank, then say \emph{nāko} again
\end{itemize}

\begin{tcolorbox}[breakable,colback=teal!4,colframe=teal!35!black,title={Before you move on}]
What does \mw{नाको} mean? (\textquotedblleft{}A checkpoint; the mouth of a road\textquotedblright{}.) Say it once more.
\end{tcolorbox}
\section[\texorpdfstring{\mw{फाटक} (phāṭak)}{फाटक (phāṭak)}]{\mw{फाटक} (phāṭak) — a big gate; a railway crossing}

\label{lesson:MW-C207-phatak}



\begin{quote}
Before the new one: say the Marwadi for a water tank, then the Marwadi for a checkpoint.
\end{quote}

\subsection*{You'll want to know: \mw{फाटक}}
\textbf{\mw{फाटक}} — \emph{phāṭak} — \textquotedblleft{}a big gate; a railway crossing\textquotedblright{}.

\begin{quote}
\mw{फाटक} \mw{बंद} \mw{है}, \mw{धीरे} \mw{चालो।}
\end{quote}

\begin{quote}
\emph{phāṭak band hai, dhīre chālo.}
\end{quote}

\begin{quote}
The gate is shut, go slowly.
\end{quote}

A masculine noun.

\begin{cousinweb}
\textbf{\mw{फाटक}} (\emph{phāṭak}) is inherited from a reconstructed \emph{*phāṭṭakka}, \textquotedblleft{}gate, door\textquotedblright{}. CDIAL gives no Marwari form in this entry, but lists Hindi \emph{phāṭak}, \textquotedblleft{}gate\textquotedblright{} (CDIAL 9069).
\end{cousinweb}

\begin{grammarlens}[title={What you've built}]
One more word for this chapter.
\end{grammarlens}

\subsection*{Your turn}
\begin{itemize}
\raggedright
  \item \emph{Say it:} \emph{phāṭak}
  \item \emph{Say it:} \emph{phāṭak}, once more
  \item \emph{Say it:} \emph{phāṭak band hai, dhīre chālo.} — The gate is shut, go slowly.
  \item \emph{Recall:} say the Marwadi for a water tank, then the Marwadi for a checkpoint, then say \emph{phāṭak} again
\end{itemize}

\begin{tcolorbox}[breakable,colback=teal!4,colframe=teal!35!black,title={Before you move on}]
What does \mw{फाटक} mean? (\textquotedblleft{}A big gate; a railway crossing\textquotedblright{}.) Say it once more.
\end{tcolorbox}
\section[\texorpdfstring{\mw{बुरज} (buraj)}{बुरज (buraj)}]{\mw{बुरज} (buraj) — a tower (of a fort)}

\label{lesson:MW-C207-buraj}



\begin{quote}
Before the new one: say the Marwadi for a checkpoint, then the Marwadi for a big gate.
\end{quote}

\subsection*{You'll want to know: \mw{बुरज}}
\textbf{\mw{बुरज}} — \emph{buraj} — \textquotedblleft{}a tower (of a fort)\textquotedblright{}.

\begin{quote}
\mw{बुरज} \mw{घणो} \mw{जूनो} \mw{है।}
\end{quote}

\begin{quote}
\emph{buraj ghaṇo jūno hai.}
\end{quote}

\begin{quote}
The tower is very old.
\end{quote}

A masculine noun.

\begin{cousinweb}
\textbf{\mw{बुरज}} (\emph{buraj}) is usually traced, through Persian, to Arabic \emph{burj}, \textquotedblleft{}a tower\textquotedblright{}.
\end{cousinweb}

\begin{grammarlens}[title={What you've built}]
One more word for this chapter.
\end{grammarlens}

\subsection*{Your turn}
\begin{itemize}
\raggedright
  \item \emph{Say it:} \emph{buraj}
  \item \emph{Say it:} \emph{buraj}, once more
  \item \emph{Say it:} \emph{buraj ghaṇo jūno hai.} — The tower is very old.
  \item \emph{Recall:} say the Marwadi for a checkpoint, then the Marwadi for a big gate, then say \emph{buraj} again
\end{itemize}

\begin{tcolorbox}[breakable,colback=teal!4,colframe=teal!35!black,title={Before you move on}]
What does \mw{बुरज} mean? (\textquotedblleft{}A tower (of a fort)\textquotedblright{}.) Say it once more.
\end{tcolorbox}
